
% \IEEEPARstart{C}{ombinatorial} optimization problems (COPs) such as the Traveling Salesman Problem (\TSP{}) and the capacitated Vehicle Routing Problem (\CVRP{}) have broad applications in logistics and network design~\cite{TSPAPP}.
% As NP-hard problems, they are typically tackled with heuristics.
% Classical solvers such as LKH-3~\cite{lkhtsp, lkhcvrp} and HGS~\cite{hgs} achieve near-optimal precision through decades of engineering, while recent Machine Learning (ML) based neural combinatorial optimization (NCO)~\cite{MLCO, RL4CO,neuopt,2optlearning} has emerged as a data-driven alternative. However, end-to-end NCO models often struggle with scalability, generalization, and a lack of convergence guarantees~\cite{RL4CO}. Learning-guided optimization (LGO)~\cite{learn2branch, learn2delegate, neurolkh} bridges this gap by integrating neural signals into classical human solvers, which often achieves state-of-the-art performance by effectively combining the strengths of both paradigms~\cite{learn2delegate,deepaco,learn2seg,huang2024contrastive}.

% Among the LGO methods, learning to guide Ant Colony Optimization (ACO)~\cite{ACO, maxmin} serves as a foundational approach for steering meta-heuristics~\cite{deepaco,gtgaco}. ACO is a population-based metaheuristic where artificial ants construct solutions via probabilistic transitions biased by pheromone updates toward promising regions. In general, ACO factorizes the transition rule by combining dynamic pheromone updates with heuristic priors. While recent methods successfully learn to initialize such heuristic priors~\cite{deepaco,gfacs} or both heuristics and pheromone matrices~\cite{gtgaco}, they all share a critical flaw: they train a model to produce guidance from instance geometry \emph{only once}, yielding static heatmaps, and then deploy this fixed signal inside a full ACO loop where pheromone dynamics evolve over hundreds of iterations.
% The model never observes pheromone during training and has no mechanism to adapt to it at inference.
% This \emph{training-inference misalignment} renders the neural signal blind to search progress: it applies identical guidance regardless of whether pheromones are near-uniform (early search) or concentrated around a local optimum (late search).
% % HeatACO~\cite{heataco} confirms that under this static paradigm the model architecture is interchangeable, suggesting that the static integration strategy itself is the limiting factor.

% Motivated by this, we propose \textbf{\DyNACO{}}, a framework that shifts from \emph{static} to \emph{dynamic} neural guidance. While prior methods train the model on a single snapshot of instance geometry and freeze the output, \DyNACO{} trains its policy on the \emph{full search trajectory}: the model repeatedly observes the evolving pheromone distribution and incumbent solution, and learns to emit guidance that is useful at every stage of the search, not just at initialization. To enable this, we design a specialized module to extract representations that summarize the current optimization status. 
% We formalize this as a \emph{semi-Markov Decision Process} in which a meta-policy periodically observes search state and emits updated edge-level guidance, and we train \DyNACO{}  with a reinforcement learning algorithm.

% Another key limitation of existing neural-guided ACO methods is scalability. To make \DyNACO{} tractable at scale, we pair the policy with a perturbation-based ACO backend on a sparse $K$-nearest-neighbor graph and a \emph{scope-restricted refinement} (SRR) mechanism. Because the resulting perturbations are highly localized, our design delivers three simultaneous benefits: (1) \emph{Stable Credit Assignment}: confining refinement to the neighborhood of the policy's actions preserves the causal link between those actions and post-refinement rewards; (2) \emph{Within-Scope Convergence}: SRR exhaustively applies 2-opt moves around the perturbed edges until no further improvement, ensuring local optimality that guarantees performance; and (3) \emph{Size-Independent Scalability}: the per-ant refinement cost is reduced from $\mathcal{O}(N^2)$ to $\mathcal{O}(M \cdot K)$, making stable training feasible at the 100,000-node scale.

% Extensive experiments on TSP and CVRP demonstrate that \DyNACO{} significantly outperforms existing neural ACO baselines under matched budgets, and surpasses established classical and neural baselines, especially on real-world large-scale instances. On TSP, neural guidance actually \emph{reduces} total runtime by 23-33\% over the unguided solver, as better-targeted perturbations accelerate local-search convergence; on CVRP, neural overhead stays below 1\%, and end-to-end wall-clock overhead remains within 1-3\% at large scales, with consistent quality gains. Finally, we conduct in-depth analyses revealing the learned guidance patterns and explaining their superiority over static priors: in particular, the policy actively counteracts ACO stagnation~\cite{stagnation} by suppressing over-reinforced edges and redirecting search toward under-explored regions. Our work underscores the necessity of aligning neural training with iterative search dynamics in learning-guided optimization.

% % \noindent\textbf{Contributions.}
% Our contributions are as follows:
% (i)~we propose \DyNACO{}, the first framework for dynamic neural guidance in ACO, formulated as a semi-MDP that lets a policy inject updated guidance at regular intervals rather than committing to a single static prediction;
% (ii)~we introduce SRR, confining local search to the perturbation neighborhood to preserve credit assignment and achieve 2-opt optimality with efficiency;
% (iii)~we show that \DyNACO{} scales to 100K nodes while reducing total runtime on TSP and adding $<$1\% overhead on CVRP, with zero-shot transfer on TSPLIB and CVRPlib confirming cross-scale and cross-distribution generalization.

\IEEEPARstart{C}{ombinatorial} optimization problems (COPs), such as the Traveling Salesman Problem (\TSP{}) and the Capacitated Vehicle Routing Problem (\CVRP{}), have broad applications in logistics, transportation, and network design~\cite{TSPAPP}. As NP-hard problems, they are typically tackled by heuristic solvers. Classical solvers such as LKH-3~\cite{lkhtsp,lkhcvrp} and HGS~\cite{hgs} achieve near-optimal performance through decades of manual algorithm engineering, while recent neural combinatorial optimization (NCO) methods~\cite{MLCO,RL4CO,neuopt,2optlearning} aim to learn heuristic rules automatically from data. Despite their promise, the practical adoption of neural routing solvers remains limited by generalization. A solver trained on one distribution should remain effective when directly applied to instances with different scales, structures, and benchmark distributions. However, recent studies show that many neural routing solvers degrade sharply under zero-shot evaluation, and some even fall behind simple construction heuristics such as nearest neighbor and random insertion~\cite{nrs_survey}.

Learning-guided optimization (LGO)~\cite{learn2branch,learn2delegate,neurolkh} offers a promising path toward more reliable neural solvers by embedding learned signals into classical heuristic frameworks. Instead of replacing the entire solver with an end-to-end neural policy, LGO improves specific algorithmic components while preserving the robustness and scalability of mature solvers~\cite{learn2delegate,deepaco,learn2seg,huang2024contrastive}. Among these methods, learning to guide Ant Colony Optimization (ACO)~\cite{ACO,maxmin} is a representative direction for neural assistance in population search~\cite{deepaco,gtgaco}. ACO constructs candidate solutions through probabilistic transitions biased by heuristic priors and dynamically updated pheromones. Its population mechanism makes it naturally suitable for exploration and adaptation, but current neural ACO methods use this structure only partially.

Most existing neural ACO methods learn static edge guidance from the input instance: some initialize the heuristic prior~\cite{deepaco,gfacs}, while others learn both heuristic and pheromone-like matrices~\cite{gtgaco}. In these designs, the neural model observes instance geometry once, and the resulting heatmap remains fixed throughout the ACO search. The solver it guides is not fixed. Pheromones concentrate, incumbent solutions change, and the population gradually shifts from broad exploration toward exploitation or stagnation. A heatmap learned from a narrow source distribution may therefore become unreliable when benchmark instances induce different search trajectories.

This mismatch raises the central question for neural-guided ACO: can neural guidance be learned as a state-dependent signal that remains useful as the ACO trajectory evolves? We answer this question with \textbf{\DyNACO{}} (Dynamic Neural Ant Colony Optimization), a dynamic base model for neural ACO\@. Instead of emitting a single heatmap before search begins, \DyNACO{} learns a policy over the search trajectory. The policy periodically observes the current pheromone distribution and incumbent solution, extracts a representation of the optimization state, and emits updated guidance over candidate edges for the next search phase. We formalize this interaction as a \emph{semi-Markov Decision Process}, where each high-level policy action influences a block of low-level ACO iterations. Neural guidance thereby becomes part of the iterative solver state rather than a fixed pre-processing signal.

\DyNACO{} resolves the state mismatch, but ACO also derives strength from population search. At a macro-step, a single dynamic guidance field gives the whole colony one learned transition bias, even though different ant groups can usefully explore different continuations around the same pheromone state. We therefore introduce \textbf{\DyNACS{}} (Dynamic Neural Ant Colony Search), which gives the dynamic colony several learned search behaviors. A shared state encoder produces behavior-conditioned trajectory distributions for different ant groups. These groups sample different perturbation trajectories, contribute to the same refined candidate pool, and update the same incumbent solution and pheromone process. A learned allocation policy assigns ant mass from the current search state, allowing the colony to change its population composition as the search evolves.

The learned guidance must also remain trainable at the scales where zero-shot generalization matters. We therefore pair \DyNACO{} and \DyNACS{} with a perturbation-based ACO environment on a sparse $K$-nearest-neighbor graph. Each policy action changes a localized part of the incumbent solution, and \emph{scope-restricted refinement} (SRR) projects the resulting trajectory back to a locally consistent solution by refining only around the affected edges. This keeps rewards tied to policy actions, avoids spending refinement effort on unrelated parts of the tour, and allows the same training and inference procedure to scale to very large routing instances.

We evaluate \DyNACS{} under a strict zero-shot protocol for both \TSP{} and \CVRP{}. A single model is trained only on 1K-node synthetic instances and is then tested directly on benchmark datasets with different scales and structures, without retraining, fine-tuning, or tuning for individual instances. \DyNACS{} consistently improves the unguided backend and prior neural baselines. Further analyses show adaptive guidance beyond short-edge preference: the policy responds to the evolving pheromone landscape, suppresses over-reinforced edges when the search begins to stagnate, and maintains complementary behaviors across the ant population.

The contributions of the paper are summarized as follows:
\begin{itemize}
\item We propose \DyNACO{}, a dynamic neural ACO framework formulated as a semi-MDP, where a policy periodically observes pheromones and incumbent solutions to update guidance throughout the ACO trajectory.
\item We introduce \DyNACS{}, which extends dynamic guidance from one state-aware transition bias to a population of behavior-conditioned trajectory distributions within the same colony.
\item We design a scalable ACO backend based on perturbations with scope-restricted refinement, preserving stable credit assignment while enabling efficient local convergence on sparse graphs up to the 100,000-node scale.
\item We evaluate generalization by training only on 1K-node instances and testing zero-shot on diverse benchmark datasets, where \DyNACS{} outperforms prior neural baselines under the same protocol.
\end{itemize}

The remainder of this paper is organized as follows. Section~\ref{sec:related_work} reviews neural combinatorial optimization, learning-guided routing solvers, neural ACO, and diversity-oriented search. Section~\ref{sec:prelim} introduces the considered routing problems and the ACO transition rule. Section~\ref{sec:method} develops the move from static neural ACO to \DyNACO{}, including the semi-MDP formulation, state representation, and trajectory-aware training. Section~\ref{sec:multi-trajectory} extends \DyNACO{} to \DyNACS{} through behavior-conditioned ant groups, shared pheromone and incumbent memory, and learned allocation. Section~\ref{sec:computational-alignment} gives the scalable search backend, including sparse candidate graphs, perturbation-based ACO, scope-restricted refinement, and complexity analysis. Section~\ref{sec:experiments} reports the experiments and analysis, including benchmark comparisons, ablations, sensitivity, efficiency, and learned guidance behavior. Finally, Section~\ref{sec:conclusion} concludes the paper and discusses future directions.
